\begin{abstract}
Running large language models efficiently today almost always means using NVIDIA GPUs. This project asks whether a fundamentally different type of AI accelerator -- Tenstorrent Wormhole -- can match or exceed GPU performance for real production models, and documents the engineering work required to find out.

I ported IBM's Granite-4.0-h (tiny and small variants) from the standard HuggingFace/PyTorch stack to Tenstorrent's Wormhole hardware, specifically a Galaxy 32 multi-chip system. The model combines recurrent and attention layers with sparse expert routing, making it a demanding and representative inference target.

The porting work surfaced three main challenges: maintaining numeric correctness when compressing model weights to fit within device memory, restructuring the model to keep data on-chip across all layers without round-tripping through the host, and working around operations that the hardware's software library does not yet support. The dominant performance bottleneck is that the host CPU must send a separate command to the device for every operation, and this communication cost -- not the arithmetic -- limits throughput. I addressed this with three optimisations. First, I recorded the full forward pass for one token once and replayed it with a single command. Second, I fused the small operations that update the model's running memory state. Third, I fused the repeated scale-and-add step used in every layer.

With trace enabled, the tiny model (4 Wormhole chips) achieves 10.65--10.69 tok/s decode throughput -- exceeding an NVIDIA A100 GPU (8.5--9.0 tok/s) by ${\approx}18\%$ and ${\approx}4\times$ faster than CPU. The small model (8 chips, trace enabled) achieves 5.84--5.89 tok/s, which is below the A100 (7.1--8.2 tok/s) by ${\approx}20\%$ due to DRAM bandwidth constraints at larger model scale, but ${\approx}9\times$ faster than CPU. These results show that non-GPU hardware can exceed A100 throughput for smaller models; the challenges and solutions are described in detail for future practitioners.
\end{abstract}
